\documentclass[dvisvgm,tikz,border=3pt]{standalone}
\usepackage{amsmath,amssymb}
\usepackage{pgfplots}
\usepackage{CJKutf8}
\pgfplotsset{compat=1.18}
\usetikzlibrary{arrows.meta,calc,positioning}

\definecolor{themebg}{HTML}{30362d}
\definecolor{themefg}{HTML}{edf4e8}
\definecolor{themegray}{HTML}{9ea897}
\definecolor{themecurve}{HTML}{7eb6ff}
\definecolor{themealert}{HTML}{ff7b7b}
\definecolor{themeamber}{HTML}{f5b942}

\pgfdeclarelayer{background}
\pgfdeclarelayer{foreground}
\pgfsetlayers{background,main,foreground}

\begin{document}
\begin{CJK*}{UTF8}{gbsn}
\pagecolor{themebg}
\begin{tikzpicture}[color=themefg, text=themefg]

\begin{axis}[
    width=13.5cm,
    height=7.2cm,
    axis lines=middle,
    axis line style={color=themefg, -{Stealth}},
    tick style={color=themefg},
    tick label style={color=themefg, font=\scriptsize},
    every axis label/.style={color=themefg},
    xmin=-0.5, xmax=11.5,
    ymin=-0.3, ymax=2.5,
    xtick={1,2,3,4,5,6,7,8,9,10},
    ytick={0,1,2},
    clip=false,
    xlabel={$n$},
    ylabel={$a_n$},
    title style={font=\footnotesize\bfseries, at={(0.5,1.06)}, anchor=south},
    title={单调有界准则：单调负责“不能回退”，有界负责“不能逃逸”，逼近上确界 $M=\sup a_n$},
]
\begin{pgfonlayer}{background}
    % 上确界 M = 2.0 虚线
    \draw[themealert, dashed, line width=1.5pt] (axis cs:0,2.0) -- (axis cs:11.0,2.0);
    \node[text=themealert, font=\scriptsize, fill=themebg, inner sep=1pt, anchor=south east] at (axis cs:11.0,2.05) {上确界 $M=\sup a_n$};

    % (M - \varepsilon, M] 带状区域
    \fill[themealert!18!themebg, opacity=0.45] (axis cs:0,1.6) rectangle (axis cs:11.0,2.0);
    \draw[themeamber, densely dashed, line width=1pt] (axis cs:0,1.6) -- (axis cs:11.0,1.6);
    \node[text=themeamber, font=\scriptsize, fill=themebg, inner sep=1pt, anchor=north east] at (axis cs:11.0,1.55) {$M-\varepsilon$};

    % 阶梯引导虚线
    \draw[themegray!40, dotted] (axis cs:4,0) -- (axis cs:4,1.625);
\end{pgfonlayer}

% 严格单调递增数列 a_n = 2 - 1.5/n
% a1 = 0.5, a2 = 1.25, a3 = 1.5, a4 = 1.625, a5 = 1.7, a6 = 1.75, a7 = 1.785, a8 = 1.812, a9 = 1.833, a10 = 1.85
\addplot[only marks, mark=*, mark size=2.8pt, fill=themecurve, draw=themefg] coordinates {
    (1, 0.5) (2, 1.25) (3, 1.5)
};
\addplot[only marks, mark=*, mark size=2.8pt, fill=themeamber, draw=themefg] coordinates {
    (4, 1.625) (5, 1.7) (6, 1.75) (7, 1.785) (8, 1.812) (9, 1.833) (10, 1.85)
};

\begin{pgfonlayer}{foreground}
    \node[text=themecurve, font=\scriptsize, fill=themebg, inner sep=0.8pt, anchor=south east] at (axis cs:1,0.55) {$a_1$};
    \node[text=themecurve, font=\scriptsize, fill=themebg, inner sep=0.8pt, anchor=south east] at (axis cs:2,1.3) {$a_2$};
    \node[text=themecurve, font=\scriptsize, fill=themebg, inner sep=0.8pt, anchor=south east] at (axis cs:3,1.55) {$a_3$};
    \node[text=themeamber, font=\scriptsize, fill=themebg, inner sep=0.8pt, anchor=north west] at (axis cs:4,1.6) {$a_N \in (M-\varepsilon, M]$};

    \node[text=themefg, font=\scriptsize, fill=themebg, inner sep=1pt, anchor=west] at (axis cs:5.2,0.6) {
        $\forall n \ge N,\ a_n \in (M-\varepsilon, M] \implies \lim_{n\to\infty}a_n = M$
    };
\end{pgfonlayer}

\end{axis}
\end{tikzpicture}
\end{CJK*}
\end{document}
